\section{Methods}

\subsection{Introduction}
From our previous works:
\begin{itemize}
\item Robustness = ability to withstand delays without affecting outage completion
\item Resilience = ability to adapt/recover/resequence under disruptions
\end{itemize}

Also:
\begin{itemize}
\item Activity durations are inherently uncertain distributions, not point values
\item Critical Path (CP) can change dynamically under uncertainty
\item Emergent work is expected and must be absorbed
\end{itemize}

So the set of developed metrics must measure:
\begin{itemize}
\item Uncertainty absorption
\item Schedule flexibility
\item Recovery capability
\item Disruption impact propagation
\end{itemize}

\subsection{Problem Formulation}

A nuclear plant outage consists of a finite set of scheduled activities, subject to precedence, resource, and operational constraints. In practice, outage performance is degraded by three main classes of uncertainty:
\begin{itemize}
\item Activity duration uncertainty
\item Emergent/unplanned work
\item Resource availability and contention
\end{itemize}

Your documents already establish that outage schedules should be evaluated not only by nominal critical path logic, but by their ability to absorb duration variability and scope growth, and by the fact that critical path structure can change under uncertainty. They also distinguish robustness from resilience: robustness is the ability to withstand delays without changing outage completion, while resilience is the ability to adapt and recover under disruption.

We therefore define a formal framework in which an outage schedule is a stochastic, constrained networked process.

A nuclear power plant outage is modeled as a stochastic resource-constrained project scheduling problem (RCPSP) subject to precedence, resource, spatial, and operational constraints. The goal is to quantify schedule robustness (resistance to disruptions) and resilience (ability to recover from disruptions).

Let:
\begin{itemize}
\item $\mathcal{A} = \{1,\dots,n\}$ be the set of planned activities
\item $\mathcal{E}$ be the set of emergent activities
\item $\mathcal{R} = \{1,\dots,m\}$ be the set of resource types
\item $\mathcal{Z} = \{1,\dots,q\}$ be the set of spatial zones
\end{itemize}

Each activity $i \in \mathcal{A} \cup \mathcal{E}$ is defined by:
\begin{itemize}
\item Planned duration $p_i$, stochastic duration $D_i$
\item Planned start time $s_i$, realized start time $S_i$
\item Resource demand $x_{ir}$
\item Zone assignment $z_i \in \mathcal{Z}$
\item Predecessor set $P(i)$
\end{itemize}

The outage completion time is:
\begin{equation}
C = \max_{i} (S_i + D_i)
\end{equation}

The baseline completion time is:
\begin{equation}
C^0 = \max_{i \in \mathcal{A}} (s_i + p_i)
\end{equation}

\subsection{Constraints}

\subsubsection{Precedence Constraints}
\begin{equation}
S_i \ge F_j \quad \forall j \in P(i)
\end{equation}

\subsubsection{Resource Constraints}
\begin{equation}
\sum_i x_{ir} \mathbf{1}\{S_i \le t < F_i\} \le A_r(t)
\end{equation}

\subsubsection{Spatial Constraints}
\begin{equation}
\sum_{i: z_i = z} \mathbf{1}\{S_i \le t < F_i\} \le K_z(t)
\end{equation}

\subsubsection{System-State Constraints}
\begin{equation}
\mathbf{1}\{S_i \le t < F_i\} \le M_k(t)
\end{equation}

\subsection{Stochastic Modeling}

\subsubsection{Activity Duration Uncertainty}
\begin{equation}
D_i \sim \mathcal{P}_i
\end{equation}

with mean $\mu_i$ and variance $\sigma_i^2$.

\subsubsection{Emergent Work}
Emergent activities arrive dynamically and are characterized by arrival time, duration, and resource requirements.

\subsubsection{Resource Availability}
\begin{equation}
A_r(t) = \bar{A}_r(t) - L_r(t)
\end{equation}

\subsection{Robustness and Resilience}

\subsubsection{Robustness}
\begin{equation}
\mathrm{Rob} = \mathbb{P}(C \le C^0)
\end{equation}

\subsubsection{Resilience}
For disruption scenario $\omega$:
\begin{equation}
\mathrm{Res}(\omega) =
\frac{C^{nr}(\omega) - C^{ar}(\omega)}
{C^{nr}(\omega) - C^0}
\end{equation}

\begin{equation}
\mathrm{Res} = \mathbb{E}[\mathrm{Res}(\omega)]
\end{equation}

\subsection{Metrics}

\subsubsection{Activity-Level Metrics}

\begin{equation}
\varepsilon_i = \frac{D_i - p_i}{p_i}
\end{equation}

\begin{equation}
AUI_i = \frac{\sigma_i}{\mu_i}
\end{equation}

\begin{equation}
TOP_i = \mathbb{P}(D_i > p_i)
\end{equation}

\begin{equation}
FCR_i = \mathbb{P}(D_i - p_i > TF_i)
\end{equation}

\subsubsection{Schedule-Level Metrics}

\begin{equation}
\Delta C = C - C^0
\end{equation}

\begin{equation}
CTV = \mathrm{Var}(C)
\end{equation}

\begin{equation}
DEP(\tau) = \mathbb{P}(C - C^0 > \tau)
\end{equation}

\subsubsection{Critical Path Metrics}

\begin{equation}
CPV = 1 - \mathbb{P}(CP(\omega) = CP^0)
\end{equation}

\subsubsection{Emergent Work Metrics}

\begin{equation}
EWR = \frac{|\mathcal{E}|}{|\mathcal{A}| + |\mathcal{E}|}
\end{equation}

\begin{equation}
EWI = C - C^{base}
\end{equation}

\begin{equation}
EAE = 1 - \frac{EWI}{\sum_{e \in \mathcal{E}} D_e}
\end{equation}

\subsubsection{Resource Metrics}

\begin{equation}
\rho_r(t) = \frac{U_r(t)}{A_r(t)}
\end{equation}

\begin{equation}
RSI_r^{max} = \max_t \rho_r(t)
\end{equation}

\begin{equation}
RIV_r = \mathrm{Var}(\rho_r(t))
\end{equation}

\subsubsection{Congestion Metrics}

\begin{equation}
ZCR_z(t) = \frac{\sum_{i:z_i=z} \mathbf{1}\{S_i \le t < F_i\}}{K_z(t)}
\end{equation}

\begin{equation}
CB = \sum_z \int_0^T \max(0, ZCR_z(t) - 1) dt
\end{equation}

\subsubsection{Schedule Stability}

\begin{equation}
SN = \frac{1}{n} \sum_i \mathbf{1}\{s_i^{(k)} \neq s_i^{(k-1)}\}
\end{equation}

\subsection{Monte Carlo Estimation}

For $k = 1,\dots,N$:
\begin{enumerate}
\item Sample activity durations $D_i^{(k)}$
\item Sample emergent work realization
\item Sample resource availability
\item Generate feasible schedule
\item Compute metrics
\end{enumerate}

\begin{equation}
\widehat{\mathrm{Rob}} = \frac{1}{N} \sum_k \mathbf{1}(C^{(k)} \le C^0)
\end{equation}

---

\section{Integration with Optimization Framework}

\subsection{RCPSP Formulation}

The outage scheduling problem is formulated as a Resource-Constrained Project Scheduling Problem (RCPSP), where the objective is not only to minimize makespan but to optimize robustness and resilience metrics.

Let $\pi$ denote a priority rule used by a Schedule Generation Scheme (SGS). The SGS maps $\pi$ into a feasible schedule $S(\pi)$.

\subsection{Multi-Objective Optimization}

We define a composite fitness function:

\begin{equation}
\mathcal{F}(\pi) =
w_1 \cdot \mathbb{E}[C(\pi)] +
w_2 \cdot CTV(\pi) +
w_3 \cdot (1 - \mathrm{Rob}(\pi)) +
w_4 \cdot (1 - \mathrm{Res}(\pi)) +
w_5 \cdot RSI^{max}(\pi) +
w_6 \cdot CB(\pi) +
w_7 \cdot SN(\pi)
\end{equation}

where $w_i$ are user-defined weights.

\subsection{Genetic Programming for Priority Rule Evolution}

Instead of predefined priority rules (e.g., minimum slack), Genetic Programming (GP) is used to evolve priority functions:

\begin{equation}
\pi(i) = f(\text{slack}_i, AUI_i, RSI_i, z_i, \text{risk}_i, \dots)
\end{equation}

The GP searches over symbolic expressions $f(\cdot)$ to optimize $\mathcal{F}(\pi)$.

\subsection{Evaluation Loop}

For each candidate priority rule $\pi$:
\begin{enumerate}
\item Generate schedule using SGS
\item Perform Monte Carlo simulation
\item Compute robustness and resilience metrics
\item Evaluate fitness $\mathcal{F}(\pi)$
\end{enumerate}

\subsection{Key Advantage}

This formulation transforms outage scheduling from a deterministic makespan minimization problem into a stochastic, multi-objective optimization problem that explicitly accounts for:
\begin{itemize}
\item uncertainty in activity durations
\item emergent work
\item resource constraints
\item congestion and spatial conflicts
\item schedule stability
\end{itemize}

This enables the discovery of adaptive scheduling policies that are better aligned with real outage execution conditions.